\documentclass[11pt,a4paper]{article}

\usepackage[utf8]{inputenc}
\usepackage[T1]{fontenc}
\usepackage[brazilian]{babel}
\usepackage{lmodern}
\usepackage{geometry}
\usepackage{microtype}
\usepackage{booktabs}
\usepackage{longtable}
\usepackage{array}
\usepackage{amsmath}
\usepackage{enumitem}
\usepackage{xcolor}
\usepackage{listings}
\usepackage{hyperref}
\usepackage{fancyhdr}

\geometry{margin=2.3cm}
\hypersetup{
  colorlinks=true,
  linkcolor=blue!50!black,
  urlcolor=blue!60!black
}
\setlist{nosep}
\pagestyle{fancy}
\setlength{\headheight}{14pt}
\fancyhf{}
\lhead{MinBias-Pythia-Geant4}
\rhead{Guia de Execução}
\cfoot{\thepage}

\definecolor{codebg}{RGB}{247,247,247}
\definecolor{codeframe}{RGB}{210,210,210}

\lstdefinestyle{shell}{
  basicstyle=\ttfamily\small,
  backgroundcolor=\color{codebg},
  frame=single,
  rulecolor=\color{codeframe},
  breaklines=true,
  columns=fullflexible,
  keepspaces=true,
  showstringspaces=false
}

\lstdefinestyle{config}{
  basicstyle=\ttfamily\small,
  backgroundcolor=\color{codebg},
  frame=single,
  rulecolor=\color{codeframe},
  breaklines=true,
  columns=fullflexible,
  keepspaces=true,
  showstringspaces=false
}

\title{\textbf{MinBias-Pythia-Geant4}\\
Guia de Configuração e Execução pela Linha de Comando}
\author{Projeto MinBias-Pythia-Geant4}
\date{\today}

\begin{document}

\maketitle
\tableofcontents
\newpage

\section{Objetivo deste guia}

Este documento descreve, de forma operacional e reproduzível, como configurar e
executar o simulador \texttt{MinBias-Pythia-Geant4} a partir da linha de comando.

O fluxo principal do programa é

\[
\text{arquivo de configuração}
\longrightarrow
\text{sobrescritas pela CLI}
\longrightarrow
\text{validação}
\longrightarrow
\text{PYTHIA}
\longrightarrow
\text{Geant4}
\longrightarrow
\text{ROOT}.
\]

O arquivo \texttt{.conf} define um perfil-base. As opções fornecidas na linha de
comando substituem, quando suportadas, os valores desse perfil apenas para a
execução corrente.

\section{Requisitos}

Antes da execução, devem estar disponíveis no ambiente:

\begin{itemize}
  \item PYTHIA 8;
  \item Geant4;
  \item ROOT;
  \item CMake;
  \item compilador C++ compatível;
  \item Python 3;
  \item Git, para desenvolvimento e auditoria do repositório.
\end{itemize}

No ambiente de desenvolvimento de referência, normalmente utiliza-se:

\begin{lstlisting}[style=shell]
conda activate ambiente_fisica
\end{lstlisting}

Entre no diretório do projeto:

\begin{lstlisting}[style=shell]
cd /caminho/para/PythiaGeantOneStage_V06
\end{lstlisting}

\section{Interface recomendada: \texttt{run.sh}}

Para uso cotidiano, a interface recomendada é o script \texttt{run.sh}.

Ele:

\begin{enumerate}
  \item verifica a disponibilidade do PYTHIA;
  \item verifica o ambiente Geant4;
  \item valida o projeto e o arquivo de configuração;
  \item configura o CMake;
  \item compila o executável em modo \texttt{Release};
  \item executa \texttt{build/pythia\_geant};
  \item repassa ao simulador as opções adicionais fornecidas na CLI.
\end{enumerate}

A forma geral é:

\begin{lstlisting}[style=shell]
./run.sh ARQUIVO.conf [opcoes]
\end{lstlisting}

Exemplo:

\begin{lstlisting}[style=shell]
./run.sh config/production.conf \
    --events 100 \
    --mu 1 \
    --threads 2 \
    --seed 9512 \
    --output outputs/meu_minbias.root
\end{lstlisting}

\section{Execução direta do binário}

Quando o projeto já está compilado, o executável também pode ser chamado
diretamente:

\begin{lstlisting}[style=shell]
./build/pythia_geant \
    --config config/production.conf \
    --events 100 \
    --mu 1 \
    --threads 2 \
    --seed 9512 \
    --output outputs/meu_minbias.root
\end{lstlisting}

A opção \texttt{--config} é obrigatória na chamada direta do executável.

Para consultar a ajuda:

\begin{lstlisting}[style=shell]
./build/pythia_geant --help
\end{lstlisting}

ou

\begin{lstlisting}[style=shell]
./build/pythia_geant -h
\end{lstlisting}

\section{Regra de precedência}

A configuração efetiva segue a ordem:

\begin{enumerate}
  \item carregamento do arquivo \texttt{.conf};
  \item aplicação das opções da linha de comando;
  \item validação da configuração final;
  \item impressão da configuração efetiva;
  \item execução, exceto quando \texttt{--dry-run} é usado.
\end{enumerate}

Assim, se \texttt{production.conf} contiver:

\begin{lstlisting}[style=config]
events = 3000
threads = 1
seed_base = 512
mean_interactions = 50.0
\end{lstlisting}

mas a execução for:

\begin{lstlisting}[style=shell]
./run.sh config/production.conf \
    --events 200 \
    --mu 2 \
    --threads 2 \
    --seed 12345 \
    --output outputs/teste.root
\end{lstlisting}

a configuração efetiva conterá, entre outros parâmetros:

\begin{lstlisting}[style=config]
events = 200
threads = 2
seed_base = 12345
interaction_mode = poisson
mean_interactions = 2
\end{lstlisting}

Os demais parâmetros permanecem com os valores definidos no arquivo-base.

\section{Opções atualmente disponíveis na linha de comando}

\begin{longtable}{>{\ttfamily}p{4.5cm} p{7.6cm} p{3.0cm}}
\toprule
\textnormal{Opção} & Significado & Exemplo \\
\midrule
\endhead

--config &
Arquivo de configuração-base. Obrigatório no executável direto. &
config/production.conf \\

--events &
Número de bunch crossings/eventos Geant4. &
1000 \\

--mu &
Média Poisson de interações por bunch crossing. Também força
\texttt{interaction\_mode = poisson}. &
50 \\

--threads &
Número de workers Geant4. &
2 \\

--seed &
Seed-base da execução. &
9512 \\

--output &
Arquivo ROOT de saída. &
outputs/run.root \\

--physics-list &
Lista de física Geant4. &
FTFP\_BERT\_ATL \\

--production-cut-mm &
Production cut padrão, em milímetros. &
1.0 \\

--generator-mode &
Modo do gerador: \texttt{pythia} ou \texttt{single\_particle}. &
pythia \\

--particle-pdg &
Código PDG no modo de partícula única. &
11 \\

--particle-kinetic-energy-gev &
Energia cinética, em GeV, no modo de partícula única. &
50 \\

--particle-eta &
Pseudorapidez da partícula única. &
0.5 \\

--particle-phi &
Ângulo azimutal da partícula única, em radianos. &
0.0 \\

--dry-run &
Valida e imprime a configuração sem executar o transporte. &
--- \\

--help, -h &
Exibe a ajuda do programa. &
--- \\

\bottomrule
\end{longtable}

\section{Configuração-base de produção}

O arquivo \texttt{config/production.conf} é o perfil de produção de referência.
Sua estrutura atual é equivalente a:

\begin{lstlisting}[style=config]
generator_mode = pythia

events = 3000
first_bcid = 0
threads = 1
seed_base = 512

interaction_mode = poisson
mean_interactions = 50.0
fixed_interactions = 1

pythia_config = pythia_minbias.cmnd
physics_list = FTFP_BERT_ATL
production_cut_mm = 1.0

beam_sigma_x_mm = 0.0
beam_sigma_y_mm = 0.0
beam_sigma_z_mm = 0.0
beam_sigma_t_ns = 0.0

max_abs_eta = 1.8

transport_neutrinos = false
generator_audit = false
check_overlaps = false
print_every = 10

output = ../outputs/minbias_mu50.root
\end{lstlisting}

\subsection{Caminhos relativos}

Os caminhos definidos dentro do arquivo \texttt{.conf} são resolvidos
relativamente ao diretório do próprio arquivo de configuração.

Por exemplo:

\begin{lstlisting}[style=config]
pythia_config = pythia_minbias.cmnd
output = ../outputs/minbias_mu50.root
\end{lstlisting}

pode ser usado mesmo quando o programa é chamado a partir de outro diretório,
desde que a estrutura relativa ao arquivo de configuração seja preservada.

\section{Parâmetros que ainda dependem do arquivo \texttt{.conf}}

Na interface atual, os parâmetros abaixo são reconhecidos pelo arquivo de
configuração, mas ainda não possuem sobrescrita própria pela CLI:

\begin{itemize}
  \item \texttt{first\_bcid};
  \item \texttt{interaction\_mode};
  \item \texttt{fixed\_interactions};
  \item \texttt{pythia\_config};
  \item \texttt{beam\_sigma\_x\_mm};
  \item \texttt{beam\_sigma\_y\_mm};
  \item \texttt{beam\_sigma\_z\_mm};
  \item \texttt{beam\_sigma\_t\_ns};
  \item \texttt{max\_abs\_eta};
  \item \texttt{transport\_neutrinos};
  \item \texttt{generator\_audit};
  \item \texttt{check\_overlaps};
  \item \texttt{print\_every}.
\end{itemize}

Portanto, opções como

\begin{lstlisting}[style=shell]
--transport-neutrinos true
\end{lstlisting}

\textbf{não existem na interface atual}. Para ativar esse comportamento, deve-se
usar um arquivo de configuração apropriado.

\section{Minimum-bias com pile-up Poisson}

O modo normal de produção minimum-bias utiliza:

\begin{lstlisting}[style=config]
interaction_mode = poisson
mean_interactions = 50.0
\end{lstlisting}

Na linha de comando, a forma recomendada é:

\begin{lstlisting}[style=shell]
--mu 50
\end{lstlisting}

A opção \texttt{--mu} atualiza a média e força o modo de interação para
\texttt{poisson}.

Exemplo de produção:

\begin{lstlisting}[style=shell]
./run.sh config/production.conf \
    --events 3000 \
    --mu 50 \
    --threads 2 \
    --seed 50050 \
    --output outputs/minbias_mu50.root
\end{lstlisting}

\section{Número fixo de interações}

Para exigir exatamente um número fixo de interações por bunch crossing, use um
arquivo \texttt{.conf} com:

\begin{lstlisting}[style=config]
interaction_mode = fixed
fixed_interactions = 5
\end{lstlisting}

Na interface atual, não existe a opção
\texttt{--fixed-interactions}; portanto, esse modo deve ser configurado pelo
arquivo \texttt{.conf}.

\section{Dry run: validação antes do transporte}

Antes de uma produção longa, recomenda-se executar exatamente o comando desejado
com \texttt{--dry-run}:

\begin{lstlisting}[style=shell]
./run.sh config/production.conf \
    --events 10000 \
    --mu 50 \
    --threads 2 \
    --seed 9512 \
    --output outputs/minbias_mu50_10k.root \
    --dry-run
\end{lstlisting}

O programa imprime a configuração efetiva e encerra antes do transporte Geant4.

Depois de conferir os parâmetros, execute novamente sem
\texttt{--dry-run}:

\begin{lstlisting}[style=shell]
./run.sh config/production.conf \
    --events 10000 \
    --mu 50 \
    --threads 2 \
    --seed 9512 \
    --output outputs/minbias_mu50_10k.root
\end{lstlisting}

Fluxo recomendado:

\[
\text{montar comando}
\rightarrow
\texttt{--dry-run}
\rightarrow
\text{conferir}
\rightarrow
\text{remover \texttt{--dry-run}}
\rightarrow
\text{produção}.
\]

\section{Exemplos minimum-bias}

\subsection{Teste pequeno, \texorpdfstring{$\mu=1$}{mu=1}}

\begin{lstlisting}[style=shell]
./run.sh config/production.conf \
    --events 100 \
    --mu 1 \
    --threads 2 \
    --seed 9512 \
    --output outputs/minbias_mu1.root
\end{lstlisting}

\subsection{Campanha com \texorpdfstring{$\mu=10$}{mu=10}}

\begin{lstlisting}[style=shell]
./run.sh config/production.conf \
    --events 1000 \
    --mu 10 \
    --threads 2 \
    --seed 10010 \
    --output outputs/minbias_mu10.root
\end{lstlisting}

\subsection{Campanha com \texorpdfstring{$\mu=50$}{mu=50}}

\begin{lstlisting}[style=shell]
./run.sh config/production.conf \
    --events 3000 \
    --mu 50 \
    --threads 2 \
    --seed 50050 \
    --output outputs/minbias_mu50.root
\end{lstlisting}

\subsection{Campanha maior}

\begin{lstlisting}[style=shell]
./run.sh config/production.conf \
    --events 10000 \
    --mu 50 \
    --threads 2 \
    --seed 9512 \
    --output outputs/minbias_mu50_10k.root
\end{lstlisting}

\section{Lista de física Geant4}

O perfil de referência usa:

\begin{lstlisting}[style=config]
physics_list = FTFP_BERT_ATL
\end{lstlisting}

Pode-se sobrescrever a lista de física:

\begin{lstlisting}[style=shell]
./run.sh config/production.conf \
    --events 100 \
    --mu 1 \
    --threads 2 \
    --physics-list FTFP_BERT \
    --output outputs/minbias_ftfp_bert.root
\end{lstlisting}

Para comparações científicas, deve-se registrar explicitamente qual lista foi
utilizada. O perfil \texttt{FTFP\_BERT\_ATL} deve permanecer como referência
enquanto não houver um estudo específico justificando outra escolha.

\section{Modo de partícula única}

O mesmo executável suporta \texttt{single\_particle}.

Exemplo para um elétron de 50 GeV:

\begin{lstlisting}[style=shell]
./run.sh config/single_particle.conf \
    --generator-mode single_particle \
    --events 100 \
    --threads 2 \
    --particle-pdg 11 \
    --particle-kinetic-energy-gev 50 \
    --particle-eta 0.5 \
    --particle-phi 0.0 \
    --output outputs/electron_50gev.root
\end{lstlisting}

As validações atuais exigem, entre outros critérios:

\begin{itemize}
  \item PDG diferente de zero;
  \item energia cinética positiva e finita;
  \item \(\lvert\eta\rvert \leq \texttt{max\_abs\_eta}\);
  \item \(\phi \in [-\pi,\pi]\).
\end{itemize}

O modo de partícula única é útil para testes de resposta e validação. O foco de
produção deste projeto permanece sendo minimum-bias com PYTHIA.

\section{Arquivo ROOT e manifesto}

A saída principal é o arquivo ROOT indicado por:

\begin{lstlisting}[style=shell]
--output outputs/minha_execucao.root
\end{lstlisting}

Antes do transporte, o programa também registra um manifesto textual ao lado da
saída, usando o nome:

\begin{lstlisting}[style=shell]
outputs/minha_execucao.root.manifest.txt
\end{lstlisting}

Esse manifesto contém a configuração normalizada usada pela execução e é
importante para rastreabilidade.

Não sobrescreva resultados importantes sem antes preservá-los ou atribuir um
novo nome ao arquivo de saída.

\section{Reprodutibilidade e número de threads}

\subsection{Estado validado até o Ciclo 9}

O Ciclo 9 demonstrou:

\begin{itemize}
  \item reprodutibilidade canônica exata entre repetições com 1 thread;
  \item reprodutibilidade canônica exata entre repetições com 2 threads;
  \item diferenças determinísticas entre uma campanha 1T e uma campanha 2T.
\end{itemize}

Portanto, \textbf{na implementação atual, o número de threads faz parte da
definição da realização estatística}. Mudar de 1 para 2 threads pode alterar os
eventos gerados, mesmo mantendo a mesma seed-base.

Consequência operacional:

\begin{quote}
Para reproduzir exatamente uma campanha existente, preserve
\texttt{--threads}, \texttt{--seed}, arquivo \texttt{.conf}, versão do código e
demais parâmetros.
\end{quote}

Os resultados completos estão documentados em:

\begin{lstlisting}[style=shell]
docs/cycle-9-performance-reproducibility/results.md
\end{lstlisting}

Na máquina de referência do Ciclo 9, duas threads produziram speedup de
aproximadamente \(1.717\times\), eficiência paralela de aproximadamente
\(85.86\%\) e redução mediana de tempo de parede de aproximadamente \(41.77\%\).

Por esse motivo, 2 threads são uma boa escolha operacional nessa máquina quando
o objetivo é throughput, desde que se compreenda a política de
reprodutibilidade descrita acima.

\section{Boas práticas para campanhas}

Para toda campanha destinada a análise científica:

\begin{enumerate}
  \item mantenha o repositório em um commit conhecido;
  \item registre o SHA do commit;
  \item registre o arquivo \texttt{.conf};
  \item registre todas as sobrescritas da CLI;
  \item use uma seed explícita;
  \item preserve o número de threads;
  \item use \texttt{--dry-run} antes da produção;
  \item use nomes de saída que identifiquem a campanha;
  \item preserve o manifesto;
  \item não descarte execuções apenas por apresentarem tempo desfavorável;
  \item para benchmarks, execute repetições independentes;
  \item não misture arquivos produzidos por configurações diferentes sem
        registrar essa diferença.
\end{enumerate}

\section{Convenção sugerida para nomes de saída}

Uma convenção simples é:

\begin{lstlisting}[style=shell]
outputs/minbias_mu<mu>_nev<eventos>_t<threads>_seed<seed>.root
\end{lstlisting}

Exemplo:

\begin{lstlisting}[style=shell]
outputs/minbias_mu50_nev10000_t2_seed9512.root
\end{lstlisting}

Isso não substitui o manifesto, mas reduz ambiguidades durante a análise.

\section{Erros comuns}

\subsection{PYTHIA não encontrado}

Se aparecer uma mensagem semelhante a:

\begin{lstlisting}[style=shell]
pythia8-config nao foi encontrado no PATH
\end{lstlisting}

ative o ambiente que contém o PYTHIA antes de executar o simulador.

\subsection{Geant4 não encontrado}

Confirme \texttt{Geant4\_DIR}, \texttt{CMAKE\_PREFIX\_PATH} ou a existência de
\texttt{geant4-config} no ambiente.

\subsection{Arquivo de configuração inexistente}

Confirme o caminho fornecido a \texttt{--config} ou ao primeiro argumento de
\texttt{run.sh}.

\subsection{Opção desconhecida}

Se o programa informar:

\begin{lstlisting}[style=shell]
Opcao desconhecida
\end{lstlisting}

consulte:

\begin{lstlisting}[style=shell]
./build/pythia_geant --help
\end{lstlisting}

e verifique se o parâmetro pretendido ainda é exclusivo do arquivo
\texttt{.conf}.

\subsection{Número de threads maior que o suportado}

Se o Geant4 não tiver sido compilado com multithreading, apenas
\texttt{--threads 1} é permitido.

\section{Receita operacional recomendada}

Para uma nova produção minimum-bias:

\begin{enumerate}
  \item escolha o arquivo-base, normalmente
        \texttt{config/production.conf};
  \item escolha número de eventos;
  \item escolha \(\mu\);
  \item escolha número de threads;
  \item escolha seed;
  \item escolha um nome de saída descritivo;
  \item execute o dry run;
  \item confira a configuração impressa;
  \item execute a produção;
  \item preserve ROOT, manifesto, commit e comando utilizado.
\end{enumerate}

Exemplo de dry run:

\begin{lstlisting}[style=shell]
./run.sh config/production.conf \
    --events 10000 \
    --mu 50 \
    --threads 2 \
    --seed 9512 \
    --output outputs/minbias_mu50_nev10000_t2_seed9512.root \
    --dry-run
\end{lstlisting}

Exemplo de produção:

\begin{lstlisting}[style=shell]
./run.sh config/production.conf \
    --events 10000 \
    --mu 50 \
    --threads 2 \
    --seed 9512 \
    --output outputs/minbias_mu50_nev10000_t2_seed9512.root
\end{lstlisting}

\section{Resumo rápido}

\begin{lstlisting}[style=shell]
# Ajuda
./build/pythia_geant --help

# Dry run
./run.sh config/production.conf \
    --events 100 \
    --mu 1 \
    --threads 2 \
    --seed 9512 \
    --output outputs/teste.root \
    --dry-run

# Producao
./run.sh config/production.conf \
    --events 100 \
    --mu 1 \
    --threads 2 \
    --seed 9512 \
    --output outputs/teste.root
\end{lstlisting}

\section{Escopo e evolução da interface}

Este guia descreve a interface efetivamente disponível no repositório no
fechamento do Ciclo 9.

Uma evolução planejada é ampliar as sobrescritas pela linha de comando para que
parâmetros hoje exclusivos do arquivo \texttt{.conf} também possam ser
controlados explicitamente pela CLI.

Até que essa extensão esteja implementada e validada, este documento distingue
deliberadamente entre:

\begin{itemize}
  \item opções que realmente existem na CLI;
  \item parâmetros que devem permanecer no arquivo \texttt{.conf}.
\end{itemize}

Essa distinção evita comandos inválidos e torna as campanhas mais auditáveis.

\section{Referência do projeto}

Repositório:

\begin{center}
\url{https://github.com/ncevidanes/MinBias-Pythia-Geant4}
\end{center}

Ao relatar resultados, registre sempre o commit do repositório e a configuração
efetiva utilizada.

\end{document}
